\chapter{Engagement Highlights} \label{highlights}
\section{Overview}
This chapter defines the Rules of Engagement (RoE) for the penetration testing activity performed as part of the Penetration Testing and Ethical Hacking exam.

\section{Purpose and Objectives}
The purpose of this penetration testing activity is to evaluate the security posture of the intentionally vulnerable virtual machine \textbf{"Corrosion: 1"}.
\\

The penetration testing activities is performed in a black-box fashion.
\\

The activity is conducted as part of the Penetration Testing and Ethical Hacking exam and aims to demonstrate practical knowledge of offensive security methodologies.
\\

The objectives of the assessment are:

\begin{itemize}
    \item Identify security vulnerabilities present within the target system.
    \item Perform reconnaissance and enumeration activities.
    \item Exploit discovered vulnerabilities where technically feasible.
    \item Demonstrate penetration testing methodologies and techniques.
    \item Document discovered vulnerabilities, exploitation procedures,
    evidence, and remediation recommendations.
\end{itemize}

The penetration test is conducted exclusively for educational purposes within an authorized laboratory environment.

\section{Scope of Testing}
\subsection{In-Scope Assets}
The following asset is explicitly authorized for testing:

\begin{itemize}
    \item \textbf{Target:} Corrosion: 1 virtual machine.
    \item \textbf{Environment:} Isolated virtual laboratory network.
\end{itemize}

All services, applications, configurations, user accounts, and files contained inside the target virtual machine are considered within scope.

\subsection{Out-of-Scope Assets}
The following assets are explicitly excluded from the scope:

\begin{itemize}
    \item The physical host machine running the virtualization software.
    \item Other virtual machines present in the laboratory environment.
    \item External networks and Internet-accessible systems.
    \item Any system not explicitly identified as the target machine.
\end{itemize}

No testing activity shall intentionally affect systems outside the authorized
scope.


\section{Authorized Testing Activities}
The following activities are authorized against the target virtual machine:

\begin{itemize}
    \item Network discovery and host identification.
    \item Port scanning and service enumeration.
    \item Vulnerability identification and assessment.
    \item Web application testing.
    \item Configuration analysis.
    \item Credential testing against exposed services.
    \item Controlled exploitation of identified vulnerabilities.
    \item Privilege escalation attempts.
    \item Post-exploitation analysis limited exclusively to the target machine.
    \item Collection of technical evidence for reporting purposes.
\end{itemize}

Since the target machine is intentionally vulnerable and isolated, exploitation activities are permitted when required to demonstrate penetration testing techniques.

\section{Prohibited Activities}
The following actions are prohibited:

\begin{itemize}
    \item Attacking systems outside the defined scope.
    \item Performing denial-of-service attacks against external systems.
    \item Modifying or damaging the host operating system.
    \item Performing security tests against third-party infrastructure.
\end{itemize}

\section{Testing Methodology}
\subsection{Methodology}
The penetration testing activities will be conducted in accordance with the \textit{Penetration Testing Execution Standard (PTES)}. 
\\

PTES provides a structured framework for planning, executing, and documenting penetration tests, ensuring that the assessment follows a consistent and methodical approach. 
\\

The methodology encompasses the relevant phases of the penetration testing process, including pre-engagement interactions, intelligence gathering, threat modelling, vulnerability analysis, exploitation, post-exploitation, and reporting \cite{ptes}. 

\subsection{Tools}
The following tools were used during the assessment:

\begin{itemize}
	\item Oracle VirtualBox \cite{virtualbox};
	\item Kali Linux \cite{kali};
    \item \texttt{nmap} \cite{nmap};
    \item \texttt{DirBuster} \cite{dirbuster};
    \item \texttt{Gobuster} \cite{gobuster};
    \item \texttt{wfuzz} \cite{wfuzz};
    \item \texttt{Nessus} \cite{nessus};
    \item \texttt{Nikto} \cite{nikto};
    \item \texttt{WafW00f} \cite{wafwoof};
    \item \texttt{Exploit Database} \cite{exploitdb};
    \item \texttt{Metasploit} \cite{metasploitMetasploitPenetration};
    \item \texttt{ssh} and \texttt{sftp} \cite{ssh};
    \item \texttt{curl} \cite{curl};
    \item \texttt{zip2john} and \texttt{john} \cite{john};
    \item \texttt{C} \cite{c}, \texttt{PHP} \cite{php} and \texttt{Python} \cite{python} programming;
    \item \texttt{netcat} \cite{netcat};
    \item \texttt{WeBaCoo} \cite{webacoo}.
\end{itemize}

\section{Data Handling}
Information collected during the penetration test shall be handled responsibly.
\\

The tester agrees to:

\begin{itemize}
    \item Use collected information only for academic evaluation purposes.
    \item Preserve only evidence required for the final report.
    \item Remove temporary testing artifacts when possible.
\end{itemize}

\section{Risk Management}
Although the target system is intentionally vulnerable, penetration testing
activities may cause service interruption, system instability, or modification
of data contained within the virtual machine.
\\

The tester is responsible for maintaining the isolation of the laboratory
environment and ensuring that testing activities remain limited to the
authorized target.
\\

Virtual machine snapshots may be used to restore the environment if necessary.

\section{Reporting and Communication}
The final penetration testing report will contain:

\begin{itemize}
    \item Scope and methodology.
    \item Tools and techniques used.
    \item Identified vulnerabilities.
    \item Exploitation procedures.
    \item Evidence of successful exploitation.
    \item Security impact analysis.
    \item Recommended mitigation strategies.
\end{itemize}

The report will be submitted to Prof. Arcangelo Castiglione as part of the
Penetration Testing and Ethical Hacking examination.

\section{Duration of Testing}

Testing activities are authorized during the exam period assigned by the
course supervisor.
\\

The tester may perform penetration testing activities against the target
virtual machine during this period, provided that all rules defined in this
document are followed.

\section{Authorization}
This section provides formal authorization for the tester to perform
penetration testing activities against the specified target system, within the scope and limitations described in this chapter.